% subsection: 階比型 \, $a_{n+1}=f(n)a_n$

  \subsection{\texorpdfstring{階比型 \, $a_{n+1}=f(n)a_n$}{階比型}}\label{節：階比型}
  \begin{bookpoint}
    掛ける値が一定でなくても,初項から各段階で掛けた値をすべて掛け合わせればよい.
  \end{bookpoint}
    等比型では同じ数を掛けた.では,$n$によって変わる式$f(n)$を掛ける場合はどうなるだろうか.
    階差型と同じように直前の項を代入してみよう.その前に,積をまとめて表す記号を導入する.
    \bookterm{sequence}{数列}$x_n$の積を
    \[
      \prod_{i=1}^n x_i
      =x_1\times x_2\times\cdots\times x_n
    \]
    と書く.$\prod$（ラージパイ）は\bookterm{product}{総積}記号である.\footnote{$x_i$は番号を付けた数の並びであればよく,特定の規則は不要である.}
    $1$から$n$までの積は\textbf{階乗}といい,$n!=1\cdot2\cdots n$と書く.
    例えば$3!=6$であり,$0!=1$と定める.
    積では,次の計算が使える.
    \begin{align*}
      &\prod_{i=1}^{n}i=n!\\
      &\prod_{i=1}^{n}k=k^n\\
      &\prod_{i=1}^{n}a_ib_i=\displaystyle\prod_{i=1}^na_i\displaystyle\prod_{i=1}^nb_i\\
      &\prod_{i=1}^{n}a_i=\displaystyle\prod_{i=1+k}^{n+k}a_{i-k}\\
      &\prod_{i=1}^{n}(i+k)=\dfrac{(n+k)!}{k!}\qquad(k\ge0)
    \end{align*}
    2行目の$k$は$i$によらない定数,最後の$k$は整数である.
    4行目は番号をずらしただけである.例えば$n=3,k=1$なら両辺とも$a_1a_2a_3$となる.
    $n\ge3$で直前の項を順に代入すると,
    \begin{align*}
      a_n&=f(n-1)a_{n-1}\\
         &=f(n-1)\{f(n-2)a_{n-2}\}=\cdots\\
         &=a_1f(1)\cdots f(n-1)=a_1\prod_{k=1}^{n-1}f(k)
    \end{align*}
    $n=2$でも$a_2=a_1f(1)$より同じ式になる.この積による表示は$f(k)=0$となる場合も成り立つ.\footnote{空積は$1$とする.}
    \noindent \bookblocklink{example-ratio-1}{problem}{answer}{【例題】}\par\nobreak
    \begin{enumerate}[label=(\arabic*),parsep=3pt,format=\bookitemlink{example-ratio-1}{problem}{answer}]
      \item $\displaystyle a_1=1,\, a_{n+1}=2na_n$
      \item $\displaystyle a_1=2,\, a_{n+1}=\frac{n+2}{n}a_n$
      \item $\displaystyle a_1=1,\, a_{n+1}=n(n+1)a_n$
    \end{enumerate}

    \noindent\bookblocklink{example-ratio-1}{hint}{problem}{【方針】}\par\nobreak
    積を直接計算するほかに,両辺を同じ形にそろえる方法もある.
    例えば(1)では,$n!=n(n-1)!$を使い,両辺を$n!$で割ってみよう.
    右辺の$n$が消え,分母に$(n-1)!$が残るため,$\dfrac{a_{n+1}}{n!}$と$\dfrac{a_n}{(n-1)!}$という同じ形ができる.
    \begin{enumerate}[label=(\arabic*),parsep=3pt,format=\bookitemlink{example-ratio-1}{hint}{problem}]
      \item 両辺を$n!$で割り,$b_n=\dfrac{a_n}{(n-1)!}$と置けば等比型になります.
      \item $n(n+1)$を一つのまとまりとして,隣の添字にずらした比を計算してみましょう.倍率$\frac{n+2}{n}$が現れます.
      \item (1)で$a_n$を割る量は$(n-1)!$でした.今回は倍率$n(n+1)$に合わせて,$n!$も組み合わせます.
    \end{enumerate}

    \noindent\bookblocklink{example-ratio-1}{answer}{problem}{【解答】}\par\nobreak
    \begin{enumerate}[label=(\arabic*),parsep=3pt,format=\bookitemlink{example-ratio-1}{answer}{problem}]
      \item $\displaystyle a_n=(n-1)!\cdot 2^{n-1}$
      \item $\displaystyle a_n=n(n+1)$
      \item $\displaystyle a_n=n!(n-1)!$
    \end{enumerate}

    \noindent\bookblocklink{example-ratio-1}{solution}{problem}{【解法】}\par\nobreak
    \begin{enumerate}[label=(\arabic*),parsep=3pt,format=\bookitemlink{example-ratio-1}{solution}{problem}]
      \item $\displaystyle a_1=1,\, a_{n+1}=2na_n$\\
            両辺を$n!$で割ると,
            \[
              \frac{a_{n+1}}{n!}=2\frac{a_n}{(n-1)!}.
            \]
            $b_n=\dfrac{a_n}{(n-1)!}$なら,$b_{n+1}$の分母は$n!$になる.添字の違いに注意しよう.
            \[
              b_{n+1}=2b_n,\quad b_1=1\quad\Longrightarrow\quad b_n=2^{n-1}.
            \]
            戻すと$a_n=(n-1)!2^{n-1}$となる.$n=1$でも$0!\cdot2^0=1$である.

      \item $\displaystyle a_1=2,\, a_{n+1}=\frac{n+2}{n}a_n$
            $n(n+1)$の隣り合う比は$\dfrac{n+2}{n}$である.与式を$(n+1)(n+2)$で割ると,
            \[
              \frac{a_{n+1}}{(n+1)(n+2)}=\frac{a_n}{n(n+1)}.
            \]
            従って$\dfrac{a_n}{n(n+1)}=\dfrac{a_1}{2}=1$であり,$a_n=n(n+1)$となる.
      \item $\displaystyle a_1=1,\, a_{n+1}=n(n+1)a_n$\\
              両辺を$n!(n+1)!$で割って,
            \begin{align*}
              &\frac{a_{n+1}}{n!(n+1)!}=\frac{a_n}{n!(n-1)!}\\
              &\therefore \quad \frac{a_n}{n!(n-1)!}=\frac{a_1}{1!(1-1)!}\\
              &a_n=n!(n-1)!
            \end{align*}
    \end{enumerate}

    何で割ればよいか分からない場合は,まず$f(1)f(2)\cdots f(n-1)$を書き,約分できるかを調べてみよう.
    例題(2)なら,$\dfrac31\cdot\dfrac42\cdots\dfrac{n+1}{n-1}$で途中の因子が消える.\footnote{\bookterm{ratio}{階比}型では,\bookterm{logarithm}{対数}によって積を和に変える方法も使える場合がある.\sectionref*{節：対数型}で扱う.}

\bookstep{掛けて倍率をそろえる}\label{解釈：階比と正規化}
    先ほどは$a_n$を適当な式で割ったが,これはその逆数を掛けたとも考えられる.\footnote{項ごとに掛ける数や割る数を変えて式をそろえる操作は,\sectionref*{節：補助数列・正規化}で扱う\bookterm{normalize}{正規化}の一例である.}
    この方法を一般に書いてみよう.$a_{n+1}=f(n)a_n$の両辺に$g(n+1)$を掛けて,$g(n)a_n$と同じ値にしたい.
    $f(k)\ne0$として,$\dfrac{g(n+1)}{g(n)}=\dfrac{1}{f(n)}$となる非零の因子を選ぶ.
    与式に$g(n+1)$を掛ければ,
    \begin{align*}
      &g(n+1)a_{n+1}=g(n)a_n=\cdots=g(1)a_1\\
      &\therefore \quad a_n=\frac{g(1)a_1}{g(n)}
    \end{align*}
    $b_n=g(n)a_n$は一定になる.例えば$f(n)=n$なら$g(n)=\dfrac{1}{(n-1)!}$なので,
    \[
      \frac{a_n}{(n-1)!}=\frac{a_1}{0!},\qquad a_n=a_1(n-1)!.
    \]
